\documentclass[11pt]{article}
\usepackage[T1]{fontenc}
\usepackage{amsmath,amssymb,geometry}
\geometry{margin=1in}
\title{A redundant tight-frame erasure counterexample\\Disproof of Conjecture 00000000952}
\author{ziangni-sys}
\date{October 8, 2026}
\begin{document}
\maketitle
\paragraph{Scope.}
The conjecture bounds the worst spectral deviation after deleting $k$
vectors from an $n$-dimensional uniform tight frame by $k/n$, with
equiangular tight frames as equality cases. We use a redundant
Parseval equiangular tight frame. Its original frame operator is exactly
the identity, so the deviation is already normalized relative to the
original tight-frame bound. No arbitrary scale factor is present.

\paragraph{The actual frame.}
In $H=\mathbb R^3$ take
\[
 v_0=\tfrac12(1,1,1)^T,\quad
 v_1=\tfrac12(1,-1,-1)^T,\quad
 v_2=\tfrac12(-1,1,-1)^T,\quad
 v_3=\tfrac12(-1,-1,1)^T.
\]
Their inner products are
\[
 \langle v_i,v_j\rangle=
 \begin{cases}3/4&i=j,\\-1/4&i\ne j.\end{cases}
\]
Thus their lengths are equal and the absolute pairwise inner products
are equal: this is an equiangular family. Directly adding the four
outer products gives
\[
 S=\sum_{i=0}^3 v_i v_i^T=I_3.
\]
Indeed each diagonal entry is $4(1/4)=1$, and each off-diagonal entry
is a sum of two $1/4$ and two $-1/4$. Hence
$\sum_i|\langle v_i,x\rangle|^2=\langle Sx,x\rangle=\|x\|^2$.
This proves that the four vectors are a uniform Parseval tight frame
in dimension $n=3$, with genuine redundancy and the ETF property.
The same real matrices define a complex ETF in $\mathbb C^3$ if desired.

\paragraph{One erasure.}
Delete $v_0$, so $k=1$. The remaining frame operator and its deviation are
\[
 S_-=\sum_{i=1}^3v_iv_i^T=I_3-v_0v_0^T,
 \qquad D=S-S_-=v_0v_0^T.
\]
The nonzero vector $v_0$ is an eigenvector of $D$, with
$Dv_0=\|v_0\|^2v_0=(3/4)v_0$. Every vector orthogonal to $v_0$ lies
in the kernel of $D$. Therefore the eigenvalues of $S_-$ are
$1/4,1,1$, and the worst spectral deviation from the original
identity is
\[
 \|D\|_{\mathrm{op}}=\frac34>\frac13=\frac{k}{n}.
\]
This already occurs for an ETF, so neither the bound nor its asserted
equality description holds as written. If uniform tight means unit-norm
tight instead, multiply every vector by $2/\sqrt3$. The tight-frame
bound becomes $4/3$, the erasure defect has norm $1$, and the relative
deviation remains $1/(4/3)=3/4$. The counterexample therefore survives
that usual normalization as well. Renormalizing after erasure or changing
the denominator to the number of frame vectors would be a different
statement from the given $n$-dimensional bound.

\paragraph{Formal verification.}
The Lean project uses complex-free real Euclidean space of dimension
three and actual continuous linear frame operators, formed from
$x\mapsto\langle v_i,x\rangle v_i$. It verifies all inner products,
the full operator identity, the remaining operator and defect, and the
eigenvector equation. Its final theorem derives
$\|D\|_{\mathrm{op}}\geq3/4>1/3$ using the genuine operator norm.
Thus the formal proof needs only a lower bound and does not assume the
claimed norm value or replace the operator with a scalar certificate.
\end{document}
